Off $\Sig$ the form $G$ is invertible, so by Lemma \claimref{sheaf.determinant-identity-blow-up-expansion}(1) $\det h = \det G \cdot (-F)$ vanishes exactly where $F$ does. Along a far-side transversal: near the wall $\beta = 1$ and $s \to -\infty$ (Lemma \claimref{sheaf.determinant-identity-blow-up-expansion}(2), $\lambda < 0$ there), so $F \to -\infty$; outside the collar $\beta = 0$ and $F = 1$. By continuity $F$ vanishes at some intermediate point, and the set $\{F = 0\}$ separates the two regimes on every transversal. Signatures: on the far-side region between $\Sig$ and the first zero of $F$, $h$ is nondegenerate ($F < 0$ there, by continuity from $F \to -\infty$ and no intervening zero), and this region's closure meets $\Sig$, where the signature is $(p{+}1,q{+}1)$ by Lemma \claimref{sheaf.wall-gram}; local constancy propagates it. Across a simple zero of $F$ the determinant $\det h$ changes sign through a simple zero, so (with respect to an auxiliary frame metric) exactly one eigenvalue of $h$ crosses zero, to first order: were two or more to vanish at the crossing point, $\det h$ would vanish there to order at least two. Which way the crossing goes is decided by the Schur complement, not by the determinant alone (a crossing to $(p{+}2, q)$ would produce the same determinant sign): off $\Sig$ the block congruence of Lemma \claimref{sheaf.determinant-identity-blow-up-expansion} gives $h \cong G \perp S$ with $S = H - \beta^2 s \begin{psmallmatrix} \nu^2 & \mu\nu \\ \mu\nu & \mu^2 \end{psmallmatrix}$, $\det S = -F$ and $\operatorname{tr} S = -\beta^2 s\,(\mu^2 + \nu^2)$. At a far-side zero of $F$ we have $2\mu\nu\beta^2 s = -1$ with $\mu\nu > 0$, so $s < 0$ and $\operatorname{tr} S > 0$: as $F$ crosses zero, $S$ passes from $\det S > 0$ with positive trace, signature $(2,0)$, to $\det S < 0$, signature $(1,1)$, while the $G$-block stays at $(p{-}1, q{+}1)$. The total therefore passes from $(p{+}1, q{+}1)$ to $(p, q{+}2)$: it is a \emph{positive} direction that degenerates. For $\mu\nu < 0$: Lemma \claimref{sheaf.wall-gram}(4) gives $(p, q{+}2)$ on the wall, which is $((p{-}1){+}1, (q{+}1){+}1)$, the ascent of the far side; and the argument above mirrors, with $s \to +\infty$ on the physical side making $F \to -\infty$ there. Finally $\mu\nu = 0$ is excluded by the wall Gram's determinant $-2\mu\nu$.
